% Práctica 2 · Plan de entregas y Product Backlog · Los Gantt-sters · MDA/DGP 2026-2027
\documentclass[11pt,a4paper]{article}

\usepackage{iftex}
\ifPDFTeX
 \usepackage[T1]{fontenc}
 \usepackage[utf8]{inputenc}
 \usepackage{tgheros}
 \renewcommand{\familydefault}{\sfdefault}
\else
 \usepackage{fontspec}
 \setmainfont{TeX Gyre Heros}
 \setmonofont{Latin Modern Mono}
\fi
\usepackage[spanish,es-noshorthands,es-tabla]{babel}
\usepackage[a4paper,top=2.9cm,bottom=2.4cm,left=2.5cm,right=2.5cm,headheight=20pt,headsep=0.8cm]{geometry}
\usepackage[table]{xcolor}
\usepackage{needspace,graphicx,longtable,array,enumitem,titlesec,fancyhdr,caption}
\usepackage{etoolbox}
\usepackage{amsmath}
\usepackage[most]{tcolorbox}
\usepackage[hidelinks]{hyperref}

% ---------- Identidad visual común ----------
\definecolor{gblue}{HTML}{1F4FB4}
\definecolor{gnavy}{HTML}{1B2236}
\definecolor{ggreen}{HTML}{2E9E5B}
\definecolor{ggray}{HTML}{5F6B7A}
\definecolor{gpale}{HTML}{8D919B}
\definecolor{glight}{HTML}{E9EFFB}
\definecolor{gline}{HTML}{B1C1E5}

\newcommand{\proyecto}{Práctica 2 · MDA}
\newcommand{\hl}[1]{\textcolor{gblue}{\textbf{#1}}}
\newcommand{\bloque}[1]{\par\medskip\noindent{\bfseries\color{gnavy}#1}\par\nopagebreak\smallskip}

\color{gnavy}
\setlength{\parindent}{0pt}
\setlength{\parskip}{6pt}
\renewcommand{\baselinestretch}{1.1}
\hypersetup{pdftitle={Práctica 2 – Plan de Entregas y Product Backlog – Los Gantt-sters},pdfauthor={Los Gantt-sters},pdfsubject={Práctica 2: Plan de Entregas del Producto y Product Backlog}}

% Títulos de sección: número en recuadro azul, texto azul y filete verde
\titleformat{\section}{\Large\bfseries\color{gblue}}
 {\colorbox{gblue}{\textcolor{white}{\thesection}}}{0.6em}{}[\vspace{-4pt}{\color{ggreen}\titlerule[0.8pt]}]
\titlespacing*{\section}{0pt}{18pt}{10pt}
\titleformat{\subsection}{\large\bfseries\color{ggreen}}{\thesubsection}{0.6em}{}
\titlespacing*{\subsection}{0pt}{10pt}{4pt}

\setlist[itemize]{label={\color{gblue}\textbullet},leftmargin=1.4em,itemsep=1pt,parsep=0pt,topsep=2pt}
\captionsetup{font={small,it},labelfont={small,it},textfont={color=ggray},labelsep=period}
\setlength{\LTpre}{6pt}\setlength{\LTpost}{8pt}
\renewcommand{\arraystretch}{1.35}
\arrayrulecolor{gline}

% ---------- Cajas ----------
% Criterios de aceptación (cabecera azul)
\newtcolorbox{criterios}[1]{enhanced,breakable,colback=glight,colframe=gblue,boxrule=0.8pt,arc=2pt,
 left=6pt,right=6pt,top=4pt,bottom=4pt,fonttitle=\small\bfseries\color{white},title={Criterios de Aceptación — #1},before upper=\small}
% Recuadro verde (releases, historia patrón, acceso a Jira)
\newtcolorbox{verde}[1]{enhanced,colback=white,colframe=ggreen,boxrule=0.8pt,arc=2pt,
 left=6pt,right=6pt,top=4pt,bottom=4pt,fonttitle=\small\bfseries\color{white},title={#1},before upper=\small}
\newcommand{\ca}[2]{\item \textbf{#1:} #2}
\setcounter{tocdepth}{2}

% Cabecera y pie
\pagestyle{fancy}\fancyhf{}
\renewcommand{\headrulewidth}{0pt}
\fancyhead[L]{\raisebox{-0.25\height}{\includegraphics[height=0.55cm]{logo.png}}\hspace{0.5em}{\small\bfseries\color{gblue}LOS GANTT-STERS}}
\fancyhead[R]{\small\color{ggray}\proyecto\ ·\ Curso 2026-2027}
\fancyfoot[C]{{\color{gpale}\rule{2.5cm}{0.4pt}}\\[2pt]\small\color{ggray}\thepage}

\begin{document}
\sloppy

% ---------- Portada (formato unificado del equipo) ----------
\begin{titlepage}
\centering
\vspace*{0.2cm}
{\fontsize{30}{34}\selectfont\bfseries\color{gblue}MDA\par}
\vspace{4pt}
{\small\bfseries\color{ggreen}Metodologías de Desarrollo Ágil\par}
\vspace{0.8cm}
\includegraphics[width=4.2cm]{logo.png}\par
\vspace{0.9cm}
{\fontsize{24}{28}\selectfont\bfseries\color{gnavy}LOS GANTT-STERS\par}
\vspace{8pt}
{\color{ggreen}\rule{\linewidth}{1.2pt}\par}
\vspace{0.8cm}
{\fontsize{20}{24}\selectfont\bfseries\color{gblue}PRÁCTICA 2\par}
\vspace{6pt}
{\bfseries\color{gnavy}PLAN DE ENTREGAS DEL PRODUCTO · PRODUCT BACKLOG\par}
\vspace{1.1cm}
{\footnotesize\bfseries\color{ggray}Integrantes\par}
\vspace{4pt}
{\color{gblue}\rule{\linewidth}{1.2pt}\par}
\vspace{10pt}
\begin{tabular}{c}
Javier Miñán Molina\\[3pt]
Germán Delgado Fernández\\[3pt]
Ana Meca Cano\\[3pt]
Jaime Montávez Ruiz\\[3pt]
Álvaro Sanabria Díaz\\[3pt]
Juan Luis Jiménez Urbano
\end{tabular}\par
\vspace{14pt}
{\small\color{gpale}Universidad de Granada · Curso 2026-2027\par}
\end{titlepage}

{\hypersetup{linkcolor=gnavy}\renewcommand{\contentsname}{\color{gblue}Índice de Contenidos}\tableofcontents}
\clearpage

\section{Introducción y contexto del proyecto}
\subsection{Origen y propósito del producto}
En el marco del proyecto de Aprendizaje-Servicio \hl{AUTOVIDA} (\textit{Software Accesible para la Autonomía y la Transición a la Vida Adulta en Estudiantes NEAE}), desarrollado en colaboración directa con el \hl{Colegio de Educación Especial (CEE) Fundación Purísima Concepción} de Granada, se planifica y desarrolla el producto software \hl{Agensible}.

La aplicación surge a partir de una necesidad pedagógica y psicológica clave transmitida por el centro educativo: \textbf{¿Y si no supieras qué pasa después?} Para el alumnado con discapacidad intelectual o del espectro autista, la incertidumbre respecto a las actividades siguientes genera nerviosismo, inseguridad, fatiga cognitiva y episodios agudos de bloqueo conductual. Anticipar estructuradamente la jornada proporciona calma, confianza y las bases necesarias para una autonomía real.

\subsection{Marco metodológico ágil y justificación del Product Owner}
Para la gestión técnica del desarrollo se adopta el marco \hl{Scrum} combinado con prácticas de ingeniería ágil (XP):

\begin{itemize}
\item \textbf{Rol del Product Owner (PO):} Siguiendo estrictamente las indicaciones del guion de la práctica, al no disponer de acceso diario al cliente real en el aula, uno de los integrantes asume el rol formal de PO actuando como la voz del cliente. Dicha responsabilidad recae en el \hl{Moderador} del grupo, \textbf{Jaime Montávez Ruiz}, quien se encargará de defender las Historias de Usuario, su granularidad y su priorización ante el profesorado evaluador.
\item \textbf{Temporización e Iteraciones:} El desarrollo se organiza en \hl{tres iteraciones oficiales} (sprints) de tres semanas de duración cada una (salvo la tercera de ajuste final antes de navidades), cumpliendo el calendario del Guion General de Prácticas 2026/2027.
\item \textbf{Inclusión de Historias Técnicas (HUT):} Para visibilizar la carga de arquitectura, configuración del backend en Supabase, despliegue en Vercel y pipeline de CI/CD exigidos en la ingeniería de software profesional, el backlog combina historias funcionales orientadas al usuario con \hl{Historias de Usuario Técnicas (HUT)} explícitas en el primer sprint.
\end{itemize}

\section{Actividad 1: Brainstorming y Product Backlog completo}
\subsection{Sesión de tormenta de ideas (Brainstorming)}
El equipo completo de \textit{Los Gantt-sters} realizó una sesión intensiva de brainstorming integrando:

\begin{enumerate}[leftmargin=1.6em,itemsep=2pt]
\item El \textbf{Pliego Técnico Oficial 26/27} de la titulación y sus anexos de requisitos.
\item Las \textbf{respuestas de la entrevista con el profesor/cliente} (roles, registro de comedor, reprografía, códigos QR, etc.).
\item Las \textbf{9 Personas y Escenarios} formalizados en la Práctica 1 previa.
\item El \textbf{material real aportado por el colegio}: fichas semanales de comandas de comedor por aulas (con dietas basal, túrmix, texturizada y cetogénica), hojas de reparto mensual de menús a familias, secuencias de tareas y recomendaciones de accesibilidad cognitiva de ARASAAC y Plena Inclusión.
\end{enumerate}

A partir de este análisis se estructuraron \textbf{7 Grandes Épicas} (6 funcionales y 1 técnica) que agrupan un total de \textbf{26 Historias de Usuario} (23 funcionales y 3 técnicas), cubriendo la totalidad de las necesidades del sistema.

\subsection{Catálogo de Grandes Épicas}
\begin{itemize}
\item \textbf{Épica Técnica (ET): Infraestructura Técnica, Arquitectura Base y DevOps.} Configuración del backend en Supabase (PostgreSQL, autenticación, almacenamiento y tiempo real), despliegue en Vercel, pipeline CI/CD en GitHub Actions y arquitectura de componentes accesibles.
\item \textbf{Épica 1 (E1): Autenticación Adaptada, Perfiles y Accesibilidad Universal.} Inicio de sesión accesible mediante foto o pictograma, gestión multidireccional de tutores y personalización sensorial/motriz.
\item \textbf{Épica 2 (E2): Agenda Personal y Estructuración Temporal (``Ahora / Después'').} Visualización estructurada de la jornada, confirmación táctil de inicio/fin y temporizador visual regresivo (\textit{Time Timer}).
\item \textbf{Épica 3 (E3): Tareas Vocacionales del Centro Escolar (PTVAL).} Comanda interactiva de comedor, reparto mensual de menús a familias, reprografía en tres fases y embajadores con códigos QR.
\item \textbf{Épica 4 (E4): Ejecución Asistida y Entorno Físico.} Avisos sensoriales de transición 5 minutos antes y conmutación de vistas diaria o semanal ampliada.
\item \textbf{Épica 5 (E5): Comunicación de Necesidades Básicas, Urgencias y Feedback.} Botón directo y permanente para necesidades fisiológicas básicas (ir al baño), alertas inmediatas de emergencia y valoración emocional de las tareas.
\item \textbf{Épica 6 (E6): Gestión de Contingencias (Plan B) y Supervisión.} Activación inmediata de rutinas alternativas (\textit{Plan B}) ante incidencias, bloqueo de tareas por prioridad y analíticas de seguimiento.
\end{itemize}

\subsection{Catálogo Exhaustivo de Historias de Usuario (Product Backlog)}
Todas las historias funcionales y técnicas han sido redactadas bajo el formato canónico estándar:

\textbf{Como} [rol / equipo] \quad \textbf{quiero} [funcionalidad / infraestructura] \quad \textbf{para} [beneficio / valor de negocio]

{\small
\begin{longtable}{|>{\centering\bfseries\arraybackslash}p{1.3cm}|>{\centering\arraybackslash}p{1.1cm}|>{\raggedright\arraybackslash}p{7.4cm}|>{\centering\arraybackslash}p{1.6cm}|>{\centering\arraybackslash}p{0.8cm}|}
\hline
\rowcolor{gblue}\textcolor{white}{\textbf{ID}} & \textcolor{white}{\textbf{Épica}} & \textcolor{white}{\textbf{Título y Resumen Canónico}} & \textcolor{white}{\textbf{MoSCoW}} & \textcolor{white}{\textbf{SP}} \\ \hline
\endhead
HUT-01 & ET & \textbf{Backend y Base de Datos en Supabase:} Como equipo dev, quiero un proyecto Supabase con base de datos relacional (PostgreSQL), autenticación y almacenamiento para disponer de un entorno persistente y seguro. & Must & 3 \\ \hline
HUT-02 & ET & \textbf{Pipeline CI/CD y Calidad:} Como gestor de calidad, quiero flujo automatizado en GitHub Actions con linters para validar calidad de código. & Must & 2 \\ \hline
HUT-03 & ET & \textbf{Arquitectura UI y Design System:} Como arquitecto, quiero motor de componentes base accesibles y servicio de voz para acelerar pantallas. & Must & 3 \\ \hline
HU-01 & E1 & \textbf{Login accesible por foto/pictograma:} Como estudiante, quiero entrar tocando mi fotografía para acceder de forma autónoma. & Must & 3 \\ \hline
HU-02 & E1 & \textbf{Login estándar supervisores:} Como tutor/admin, quiero iniciar sesión con email/password para gestionar perfiles con seguridad. & Must & 2 \\ \hline
HU-03 & E1 & \textbf{Gestión de perfiles y tutores:} Como admin, quiero dar de alta estudiantes vinculando su tutor docente y un tutor familiar para coordinar rutinas. & Must & 3 \\ \hline
HU-04 & E1 & \textbf{Personalización estética (ARASAAC):} Como tutor, quiero ajustar tamaño de botones, colores y pictogramas para evitar sobreestimulación. & Must & 3 \\ \hline
HU-05 & E1 & \textbf{Interacción con pulsador (switch):} Como estudiantecon movilidad reducida, quiero barrido visual y pulsador para usar la app autónomamente. & Could & 5 \\ \hline
HU-06 & E2 & \textbf{Agenda diaria (``Ahora / Después''):} Como estudiante, quiero ver qué hago ahora y qué viene después para anticipar el día y reducir la ansiedad. & Must & 5 \\ \hline
HU-07 & E2 & \textbf{Creación de tareas personales:} Como tutor o familia, quiero programar tareas del hogar o escolares con horario y pictogramas. & Should & 3 \\ \hline
HU-08 & E2 & \textbf{Confirmación de inicio y fin:} Como estudiante, quiero botones gigantes para marcar cuándo empiezo y cuándo acabo una tarea. & Must & 2 \\ \hline
HU-09 & E2 & \textbf{Temporizador visual (Time Timer):} Como estudianteno lector de horas, quiero ver un disco visual decreciente para entender el tiempo restante. & Must & 3 \\ \hline
HU-10 & E2 & \textbf{Desglose de tareas con pasos multimedia:} Como tutor, quiero secuenciar tareas con vídeos o audios para guiar procedimientos complejos. & Should & 5 \\ \hline
HU-11 & E3 & \textbf{Comanda de comedor interactiva:} Como estudianteencargado, quiero registrar las ausencias y tipos de dieta (túrmix, basal, etc.) por aulas. & Should & 5 \\ \hline
HU-12 & E3 & \textbf{Envío de parte a cocina:} Como estudianteo tutor, quiero enviar el resumen de menús a cocina con un toque para evitar errores en las comidas. & Should & 3 \\ \hline
HU-13 & E3 & \textbf{Reparto de menús a familias:} Como estudiantede reparto, quiero una lista con fotos por aula para marcar a quién he entregado el menú en papel. & Should & 3 \\ \hline
HU-14 & E3 & \textbf{Reprografía en 3 fases:} Como estudiantede reprografía, quiero un flujo guiado (recogida, taller, entrega) para cumplir encargos del centro. & Should & 5 \\ \hline
HU-15 & E3 & \textbf{Embajador con QR y vídeo:} Como estudianteguía, quiero escanear códigos QR en las salas del colegio para reproducir vídeos explicativos en visitas. & Should & 5 \\ \hline
HU-16 & E4 & \textbf{Notificaciones anticipatorias:} Como estudiante, quiero un aviso suave 5 minutos antes de la siguiente tarea para prepararme para la transición. & Must & 2 \\ \hline
HU-17 & E4 & \textbf{Selector de vista día/semana:} Como tutor o estudianteavanzado, quiero alternar entre vista diaria o semanal según capacidad cognitiva. & Could & 2 \\ \hline
HU-18 & E5 & \textbf{Botón permanente de necesidades:} Como estudianteno verbal, quiero un botón fijo para pedir ir al baño, agua o descanso de inmediato. & Must & 3 \\ \hline
HU-19 & E5 & \textbf{Alerta de emergencia al tutor:} Como estudianteen situación de crisis, quiero un botón de auxilio de un toque que emita alarma al tutor. & Should & 3 \\ \hline
HU-21 & E5 & \textbf{Feedback emocional con caritas:} Como estudiante, quiero indicar cómo me he sentido al terminar con 3 caritas para evaluar la actividad. & Could & 2 \\ \hline
HU-22 & E6 & \textbf{Conmutación a ``Plan B'':} Como tutor o estudianteante un imprevisto, quiero activar una rutina alternativa predefinida evitando bloqueos. & Should & 5 \\ \hline
HU-23 & E6 & \textbf{Bloqueo de tareas por prioridad:} Como tutor, quiero bloquear la cancelación de tareas prioritarias para preservar la disciplina escolar. & Should & 3 \\ \hline
HU-25 & E6 & \textbf{Historial y analítica semanal:} Como tutor o familia, quiero ver el porcentaje de cumplimiento semanal para valorar la evolución del estudiante. & Could & 3 \\ \hline
\end{longtable}}

\begin{verde}{Nota sobre la numeración}
Los identificadores no son correlativos: las historias \textbf{HU-20} (chat entre tutor y estudiante) y \textbf{HU-24} (registro de asistencia del profesorado) se retiraron del Product Backlog por decisión del equipo. Se mantiene la numeración original para no alterar las referencias entre documentos.
\end{verde}

\needspace{12\baselineskip}
\section{Actividad 2: Priorización del Product Backlog}
\subsection{Criterios de priorización del Product Owner}
El Product Owner (\textbf{Jaime Montávez Ruiz}) estableció la jerarquía de prioridades sobre la base de cuatro criterios pedagógicos y técnicos:

\begin{enumerate}[leftmargin=1.6em,itemsep=2pt]
\item \textbf{Fundamentos técnicos y de entorno:} Las historias técnicas (HUT) son indispensables para habilitar la persistencia de datos, la integración continua y el motor accesible desde el día uno.
\item \textbf{Mitigación inmediata del estrés cognitivo:} Lo primero que debe resolver la herramienta es el acceso visual independiente y la visión de ``Ahora / Después''.
\item \textbf{Seguridad y necesidades fisiológicas primarias:} El botón directo de petición de baño o agua es de máxima prioridad humanitaria y operativa en el colegio.
\item \textbf{Actividades vocacionales del centro (PTVAL):} Las tareas de comedor y reprografía sustentan el aprendizaje laboral de los estudiantes y se planifican para el segundo ciclo.
\end{enumerate}

\subsection{Clasificación MoSCoW consolidada}
\begin{itemize}
\item \textbf{Must Have (12 historias - 34 SP):} HUT-01, HUT-02, HUT-03, HU-01, HU-02, HU-03, HU-04, HU-06, HU-08, HU-09, HU-16, HU-18. Conforman el núcleo del primer sprint.
\item \textbf{Should Have (10 historias - 41 SP):} HU-07, HU-10, HU-11, HU-12, HU-13, HU-14, HU-15, HU-19, HU-22, HU-23. Funcionalidades vocacionales y de resiliencia escolar.
\item \textbf{Could Have (4 historias - 11 SP):} HU-05, HU-17, HU-21, HU-25. Mejoras avanzadas de interacción motriz y analítica de datos.
\item \textbf{Won't Have (2 historias):} El equipo asume el desarrollo de las 26 historias en los tres sprints programados. Las historias HU-20 y HU-24 fueron retiradas del backlog y no se consideran Won't Have.
\end{itemize}

\section{Actividad 3: Estimación del esfuerzo (Planning Poker)}
\subsection{Técnica y escala de estimación}
El equipo empleó la técnica de \hl{Planning Poker} utilizando la escala de \textbf{Fibonacci modificada}:

\begin{center}Escala: $1,\ 2,\ 3,\ 5,\ 8,\ 13,\ 20$ Story Points (SP)\end{center}

\subsection{Historia de referencia (``Golden Story'') y reajuste por HUT}
\begin{verde}{Historia Patrón de Calibración: HU-02}
\textbf{HU-02: Inicio de sesión estándar para supervisores (email/password).} Se adopta como referencia base de 2 SP al ser una tarea de dificultad técnica conocida (formulario de entrada, autenticación con Supabase Auth y verificación del token). Al extraerse la infraestructura base del servidor y base de datos a \textbf{HUT-01} (3 SP) y el sistema de diseño accesible a \textbf{HUT-03} (3 SP), las historias \textbf{HU-03} y \textbf{HU-04} se ajustan a \textbf{3 SP cada una}, reflejando su lógica pura de negocio.
\end{verde}
La suma total de esfuerzo para las 26 historias asciende a exactamente \hl{86 Story Points}.

\section{Actividad 4: Estimación de la velocidad del equipo de desarrollo}
\subsection{Capacidad del equipo y dedicación horaria}
El dimensionamiento del equipo se deduce de los parámetros formales del curso:

\begin{itemize}
\item \textbf{Número de miembros ($N$):} 6 ingenieros de software.
\item \textbf{Duración del sprint ($T$):} 3 semanas de trabajo por iteración (15 días lectivos).
\item \textbf{Dedicación individual semanal:} 2 horas presenciales de clase de prácticas + 5.5 horas de desarrollo y coordinación no presencial = \textbf{7.5 horas/semana por persona}.
\end{itemize}
\[\text{Horas brutas de equipo por Sprint} = 6 \text{ personas} \times 7.5 \text{ h/sem} \times 3 \text{ sem} = \textbf{135}\text{ horas}\]

\subsection{Factor de enfoque (Focus Factor) y horas ideales}
Aplicando un \textbf{Factor de Enfoque} ($FF$) del \hl{65\,\%} ($0.65$) para descontar reuniones, documentación formal y clases:
\[\text{Capacidad neta de desarrollo (Horas Ideales)} = 135 \text{ h} \times 0.65 = \textbf{87.75}\text{ horas ideales}\]

\subsection{Deducción de la velocidad nominal (V)}
Tomando la ratio de \textbf{2.8 horas ideales por Story Point}:
\[\text{Velocidad estimada por Sprint } (V) = \frac{87.75 \text{ horas ideales}}{2.8 \text{ h / SP}} \approx \textbf{31}\text{ Story Points / Sprint}\]
El equipo programa absorber 34 SP en el Sprint 1 (arranque de infraestructura), 29 SP en el Sprint 2 y 23 SP en el Sprint 3, totalizando exactamente los \hl{86 Story Points} del backlog. La carga de los sprints 2 y 3 queda por debajo de la velocidad estimada, lo que deja margen para imprevistos y para el refinamiento del backlog.

\section{Actividad 5: Plan de entregas del producto (Release Plan)}
\subsection{Distribución en tres entregas incrementales}
El plan de entregas desacopla el producto en tres incrementos potencialmente utilizables por el colegio:

\begin{verde}{Release v0.1 — Entrega 1 (Sprint 1: 22 Octubre – 12 Noviembre 2026) · 34 SP}
\textbf{Objetivo:} Despliegue de la arquitectura base, backend en Supabase, pipeline CI/CD y el \textit{Producto Mínimo Viable (MVP)} de autonomía (login visual, agenda ``Ahora / Después'', temporizador visual y botón de socorro). \textbf{Historias incluidas (12 historias):} HUT-01, HUT-02, HUT-03, HU-01, HU-02, HU-03, HU-04, HU-06, HU-08, HU-09, HU-16, HU-18. \textbf{Total acumulado:} 34 SP.
\end{verde}
\begin{verde}{Release v0.2 — Entrega 2 (Sprint 2: 12 Noviembre – 3 Diciembre 2026) · 29 SP}
\textbf{Objetivo:} Operatividad vocacional del centro (PTVAL). Digitalización del control de comedor con dietas reales, reparto de menús a familias, reprografía guiada y embajadores con códigos QR. \textbf{Historias incluidas (7 historias):} HU-07, HU-10, HU-11, HU-12, HU-13, HU-14, HU-15. \textbf{Total acumulado:} 63 SP (34 + 29).
\end{verde}
\begin{verde}{Release v1.0 — Entrega 3 (Sprint 3: 3 Diciembre – 21 Diciembre 2026) · 23 SP}
\textbf{Objetivo:} Sistema integral, resiliencia y accesibilidad avanzada. Conmutación a ``Plan B'' ante imprevistos, alertas de emergencia, valoración emocional, pulsadores físicos y métricas. \textbf{Historias incluidas (7 historias):} HU-05, HU-17, HU-19, HU-21, HU-22, HU-23, HU-25. \textbf{Total acumulado:} 86 SP (100\,\% del Backlog).
\end{verde}

\section{Actividad 6: Especificación detallada y Criterios de Aceptación (Iteración 1)}
Las 12 historias del primer sprint cuentan con criterios de aceptación formales en sintaxis BDD (\textit{Behavior-Driven Development}):

\subsection{HUT-01: Configuración del backend y la base de datos en Supabase (3 SP)}
\begin{criterios}{HUT-01}
\begin{itemize}
\ca{CA-T01.1 (Proyecto Supabase operativo)}{Dado que un desarrollador clona el repositorio y configura las variables de entorno, cuando arranca la aplicación, entonces se conecta al proyecto Supabase (PostgreSQL, autenticación y almacenamiento) sin errores.}
\ca{CA-T01.2 (Migraciones automáticas)}{Dado que la base de datos se inicializa por primera vez, cuando se aplican las migraciones, entonces se crean el esquema y los datos semilla de prueba.}
\ca{CA-T01.3 (Variables de entorno)}{Dado que se configuran credenciales sensibles, cuando se despliega, entonces se aíslan en archivos \texttt{.env} ignorados por el control de versiones y en las variables de entorno de Vercel.}
\end{itemize}
\end{criterios}

\subsection{HUT-02: Pipeline de Integración Continua (CI/CD) y calidad de código (2 SP)}
\begin{criterios}{HUT-02}
\begin{itemize}
\ca{CA-T02.1 (Workflow automático)}{Dado que se abre una \textit{pull request}, cuando se sube el código, entonces GitHub Actions compila el proyecto y ejecuta los tests unitarios automáticamente.}
\ca{CA-T02.2 (Linter y formato)}{Dado que se evalúa el código, cuando se corre el pipeline, entonces se valida el formato con ESLint/Prettier y no se permite fusión si existen advertencias graves.}
\ca{CA-T02.3 (Despliegue staging)}{Dado que una rama principal es actualizada, cuando pasa los tests, entonces se genera automáticamente una versión desplegada en Vercel (entorno de pruebas).}
\end{itemize}
\end{criterios}

\subsection{HUT-03: Arquitectura base del cliente web y motor accesible (3 SP)}
\begin{criterios}{HUT-03}
\begin{itemize}
\ca{CA-T03.1 (Setup cliente)}{Dado que se inicializa el proyecto frontend, cuando se arranca la aplicación, entonces cuenta con estructura modular, router configurado y gestión de estado reactivo.}
\ca{CA-T03.2 (Componentes base accesibles)}{Dado que se maqueta una pantalla, cuando se usan los botones y contenedores base, entonces incorporan resaltado de foco cromático y áreas táctiles mínimas ($\geq 48 \times 48$ px).}
\ca{CA-T03.3 (Servicio de audio)}{Dado que se requiere accesibilidad sonora, cuando se invoca el servicio de audio, entonces reproduce locuciones y sonidos amigables de confirmación.}
\end{itemize}
\end{criterios}

\subsection{HU-01: Inicio de sesión accesible mediante foto o pictograma (3 SP)}
\begin{criterios}{HU-01}
\begin{itemize}
\ca{CA-01.1 (Cuadrícula visual)}{Dado que el estudianteabre la app, cuando accede a bienvenida, entonces visualiza las fotos de rostro del aula en tarjetas táctiles de mínimo $120 \times 120$ px con nombre en lectura fácil.}
\ca{CA-01.2 (Confirmación amigable)}{Dado que el estudiantetoca su fotografía, cuando se valida el toque, entonces se emite un sonido alegre y accede a su agenda sin requerir contraseña escrita.}
\ca{CA-01.3 (Tableta personal)}{Dado que un dispositivo está asignado a un único estudiante, cuando se inicia la aplicación, entonces salta directamente a su sesión personal.}
\end{itemize}
\end{criterios}

\subsection{HU-02: Inicio de sesión estándar para supervisores (2 SP)}
\begin{criterios}{HU-02}
\begin{itemize}
\ca{CA-02.1 (Acceso restringido)}{Dado que un tutor pulsa el icono de acceso docente en la esquina superior, cuando se abre el modal, entonces se presenta un formulario clásico de email y contraseña.}
\ca{CA-02.2 (Validación segura)}{Dado que se introducen credenciales válidas, cuando se envía el formulario, entonces el sistema genera token JWT seguro y redirige al panel según rol.}
\ca{CA-02.3 (Gestión de fallos)}{Dado que se introducen datos erróneos, cuando se procesa la petición, entonces se indica qué campo ha fallado: por ejemplo, si la contraseña es incorrecta, se muestra el error junto al campo de la contraseña, y no solo un mensaje de error general.}
\end{itemize}
\end{criterios}

\subsection{HU-03: Gestión de perfiles y tutores multidireccionales (3 SP)}
\begin{criterios}{HU-03}
\begin{itemize}
\ca{CA-03.1 (Alta de estudiante)}{Dado que el administrador crea un estudiante, cuando rellena la ficha, entonces puede asignar nombre, aula y fotografía.}
\ca{CA-03.2 (Vinculación de tutores)}{Dado que se configura el estudiante, cuando se asignan supervisores, entonces puede asociarse a su tutor de aula y a un único tutor familiar externo.}
\ca{CA-03.3 (Segregación de permisos)}{Dado que un tutor familiar inicia sesión, cuando accede a la agenda, entonces solo puede editar rutinas domésticas sin alterar tareas escolares.}
\end{itemize}
\end{criterios}

\subsection{HU-04: Personalización estética y sensorial de la interfaz (3 SP)}
\begin{criterios}{HU-04}
\begin{itemize}
\ca{CA-04.1 (Ajustes de visualización)}{Dado que el tutor entra a configuración, cuando selecciona opciones, entonces puede escoger entre 3 tamaños de botón y 3 paletas (oscuro, claro y pastel).}
\ca{CA-04.2 (Repositorio ARASAAC)}{Dado que la app muestra pictogramas, cuando se cargan los recursos, entonces se emplean exclusivamente iconos homologados de ARASAAC con alto contraste.}
\ca{CA-04.3 (Persistencia reactiva)}{Dado que el tutor guarda cambios de personalización, cuando el estudiantetiene la sesión abierta, entonces la interfaz se actualiza de inmediato.}
\end{itemize}
\end{criterios}

\subsection{HU-06: Vista de agenda estructurada del día (``Hecho - Ahora - Después'') (5 SP)}
\begin{criterios}{HU-06}
\begin{itemize}
\ca{CA-06.1 (Disposición en 3 bloques)}{Dado que el estudianteobserva su jornada, cuando mira la pantalla, entonces ve tres columnas: HECHO (última completada atenuada), AHORA (ocupando el 60\,\% central) y DESPUÉS (siguiente actividad con hora e icono).}
\ca{CA-06.2 (Avance automático)}{Dado que el estudianteconfirma el fin de la tarea en curso, cuando se registra, entonces la tarea actual pasa a HECHO y la de DESPUÉS pasa suavemente a AHORA.}
\ca{CA-06.3 (Foco atencional)}{Dado que la pantalla está activa, cuando se renderiza, entonces no se muestran barras de desplazamiento ni elementos que distraigan la atención.}
\end{itemize}
\end{criterios}

\subsection{HU-08: Confirmación accesible de inicio y fin de tarea (2 SP)}
\begin{criterios}{HU-08}
\begin{itemize}
\ca{CA-08.1 (Botón EMPEZAR)}{Dado que una tarea está en el bloque AHORA, cuando el estudiantepulsa el botón verde EMPEZAR, entonces pasa a estado ``En curso'' y activa el temporizador.}
\ca{CA-08.2 (Botón TERMINAR)}{Dado que la tarea está en curso, cuando se pulsa TERMINAR, entonces se muestra una animación de felicitación y se guarda la marca temporal exacta.}
\ca{CA-08.3 (Seguridad antipulsación)}{Dado que el estudiantepuede pulsar por error, cuando interactúa, entonces se requiere una pulsación sostenida de 1.5 segundos o confirmación simple.}
\end{itemize}
\end{criterios}

\subsection{HU-09: Temporizador visual accesible integrado (Time Timer) (3 SP)}
\begin{criterios}{HU-09}
\begin{itemize}
\ca{CA-09.1 (Disco circular regresivo)}{Dado que una tarea de duración conocida inicia, cuando corre el tiempo, entonces un disco de color vivo mengua proporcionalmente en tiempo real.}
\ca{CA-09.2 (Rótulo opcional)}{Dado que el perfil admite dígitos, cuando se dibuja el temporizador, entonces se incluye debajo un rótulo en lectura fácil (``Quedan 10 minutos'').}
\ca{CA-09.3 (Alarma visual suave)}{Dado que el disco llega a cero, cuando expira el tiempo, entonces el borde parpadea suavemente emitiendo una campana amistosa.}
\end{itemize}
\end{criterios}

\subsection{HU-16: Notificaciones sensoriales anticipatorias (2 SP)}
\begin{criterios}{HU-16}
\begin{itemize}
\ca{CA-16.1 (Disparo a los 5 minutos)}{Dado que faltan 5 minutos para la siguiente tarea, cuando se alcanza ese hito, entonces la tablet emite una vibración leve y un sonido alegre.}
\ca{CA-16.2 (Aviso flotante claro)}{Dado que se activa el aviso, cuando el estudiantemira la tablet, entonces aparece una tarjeta emergente con el pictograma de la tarea que vendrá.}
\ca{CA-16.3 (Cierre no bloqueante)}{Dado que se muestra el aviso, cuando el estudiantelo toca o pasan 15 segundos, entonces se oculta sin interrumpir la actividad actual.}
\end{itemize}
\end{criterios}

\subsection{HU-18: Botón directo permanente de necesidades fisiológicas básicas (3 SP)}
\begin{criterios}{HU-18}
\begin{itemize}
\ca{CA-18.1 (Ubicación fija)}{Dado que el estudiante navega por la app, cuando busca el botón, entonces se ubica siempre en la misma esquina inferior izquierda con icono inequívoco.}
\ca{CA-18.2 (Menú de 3 toques)}{Dado que se pulsa el botón, cuando se abre la ventana, entonces se muestran tres botones gigantes: ``Quiero ir al baño'', ``Quiero agua'', ``Necesito descanso''.}
\ca{CA-18.3 (Aviso inmediato al tutor)}{Dado que se pulsa ``Baño'', cuando se confirma, entonces el tutor escolar recibe una alarma sonora prioritaria instantánea en su móvil.}
\end{itemize}
\end{criterios}

\section{Desglose técnico en tareas de ingeniería por sprint}
Para organizar la distribución técnica del trabajo entre los 6 integrantes, cada Historia se desglosa en subtareas de ingeniería según su especialidad:

\subsection{Sprint 1 (Iteración 1) — Arquitectura Base y Núcleo de Autonomía (34 SP)}
\begin{itemize}
\item \textbf{HUT-01 (Supabase y BD):} T-T01.1 [Backend] Proyecto Supabase y conexión del cliente (Highest) · T-T01.2 [BD] Esquemas y migraciones en PostgreSQL (Highest) · T-T01.3 [Seguridad] Variables de entorno (.env) (High).
\item \textbf{HUT-02 (CI/CD):} T-T02.1 [CI] Workflow de GitHub Actions con build y tests (High) · T-T02.2 [QA] Reglas de ESLint, Prettier y a11y (High) · T-T02.3 [DevOps] Despliegue en Vercel (Medium).
\item \textbf{HUT-03 (UI Accesible):} T-T03.1 [Arquitectura] Setup cliente, router y estado reactivo (Highest) · T-T03.2 [UI] Design System con contrastes y foco grueso (Highest) · T-T03.3 [Audio] Wrapper de síntesis de voz (High).
\item \textbf{HU-01 (Login foto):} T-01.1 [UX/Front] Cuadrícula táctil accesible $\geq 120 \times 120$ px (Highest) · T-01.2 [Back] Endpoint de login visual sin password (Highest) · T-01.3 [Audio] Feedback acústico de confirmación (High) · T-01.4 [Testing] Pruebas de usabilidad cognitiva (Medium).
\item \textbf{HU-02 (Login tutor):} T-02.1 [Front] Formulario email/password validado con mensaje de error en el campo que falla (Highest) · T-02.2 [Back] Supabase Auth (JWT) y políticas de acceso por rol (Highest) · T-02.3 [Testing] Pruebas de seguridad y expiración (High).
\item \textbf{HU-03 (Perfiles):} T-03.1 [BD] Modelos relacionales estudiantes y tutores (Highest) · T-03.2 [Back] CRUD de estudiantes con foto (Highest) · T-03.3 [Front] Consola de asignación de tutor de aula y tutor familiar (High) · T-03.4 [Back] Políticas de segregación escolar/doméstica (High).
\item \textbf{HU-04 (ARASAAC):} T-04.1 [CSS] Temas de alto contraste y pastel (High) · T-04.2 [API] Catálogo oficial ARASAAC (Highest) · T-04.3 [Front] Selector de tamaños de botón y tipografías (High) · T-04.4 [QA] Auditoría de contraste WCAG 2.1 AA/AAA (High).
\item \textbf{HU-06 (Agenda Ahora/Después):} T-06.1 [Front] Componente de 3 bloques Hecho-Ahora-Después (Highest) · T-06.2 [Front] Lógica de avance reactivo (Highest) · T-06.3 [Back] Endpoint de secuencia cronológica diaria (Highest) · T-06.4 [UX] Supresión de barras de scroll (High).
\item \textbf{HU-08 (Confirmación):} T-08.1 [Front] Botones gigantes EMPEZAR y TERMINAR (Highest) · T-08.2 [UX] Prevención de toques accidentales 1.5s (High) · T-08.3 [Back] Registro de marcas de tiempo reales (High).
\item \textbf{HU-09 (Time Timer):} T-09.1 [Front/SVG] Disco circular regresivo animado (Highest) · T-09.2 [Front] Rótulo opcional en lectura fácil (Medium) · T-09.3 [Audio] Alarma no estridente al expirar tiempo (High).
\item \textbf{HU-16 (Notificación 5 min):} T-16.1 [Front] Servicio temporizador de 5 min previos (High) · T-16.2 [UX] Ventana flotante no intrusiva con vibración (High) · T-16.3 [QA] Pruebas de precisión del disparador (Medium).
\item \textbf{HU-18 (Botón necesidades):} T-18.1 [Front] Botón flotante fijado en todas las vistas (Highest) · T-18.2 [UX] Modal con 3 opciones gigantes ARASAAC (Highest) · T-18.3 [Back/Realtime] Canal en tiempo real (Supabase Realtime) de alerta sonora al tutor (Highest).
\end{itemize}

\subsection{Sprint 2 (Iteración 2) — Tareas Vocacionales del Centro (29 SP)}
\begin{itemize}
\item \textbf{HU-07 (Tareas personales):} T-07.1 [Front] Formulario con buscador ARASAAC (High) · T-07.2 [Back] API de tareas recurrentes semanales (High) · T-07.3 [Lógica] Permisos tutores vs familia (Medium).
\item \textbf{HU-10 (Pasos multimedia):} T-10.1 [BD] Modelo de subtareas ordenadas (High) · T-10.2 [Front] Reproductor paso a paso con vídeo/audio (High) · T-10.3 [Back] Almacenamiento multimedia comprimido (Medium).
\item \textbf{HU-11 (Comanda comedor):} T-11.1 [Front] Cuadrícula táctil por clases reales (High) · T-11.2 [UX] Marcado de dietas basal, túrmix, texturizada (High) · T-11.3 [Back] Registro diario de comensales (High).
\item \textbf{HU-12 (Parte a cocina):} T-12.1 [Back] Algoritmo de consolidación de menús por dieta (High) · T-12.2 [Front] Vista de recepción para cocina con exportación (Medium) · T-12.3 [UX] Animación de confirmación para el estudiante(Medium).
\item \textbf{HU-13 (Menús a familias):} T-13.1 [Front] Checklist con fotos por aula (Medium) · T-13.2 [Back] Historial mensual de entregas a familias (Medium).
\item \textbf{HU-14 (Reprografía):} T-14.1 [Back] Máquina de estados: recogida-taller-reparto (High) · T-14.2 [Front] Interfaz interactiva en 3 fases (High) · T-14.3 [Front] Formulario de encargo docente (High).
\item \textbf{HU-15 (Guía con QR):} T-15.1 [Hardware] Lector de cámara QR integrado (High) · T-15.2 [Media] Asociación de QR a salas y vídeo explicativo (High) · T-15.3 [Admin] Generador imprimible de códigos QR de salas (Medium).
\end{itemize}

\subsection{Sprint 3 (Iteración 3) — Resiliencia y Accesibilidad Avanzada (23 SP)}
\begin{itemize}
\item \textbf{HU-05 (Pulsadores):} T-05.1 [Accesibilidad] Motor de barrido visual con foco cromático (Medium) · T-05.2 [Hardware] Captura de pulsadores externos USB/switch (Medium) · T-05.3 [QA] Navegación completa con pulsador único (Low).
\item \textbf{HU-17 (Vista semanal):} T-17.1 [Front] Cuadrícula semanal accesible (Low) · T-17.2 [Front] Conmutador de vista en perfil de estudiante(Low).
\item \textbf{HU-19 (Alerta emergencia):} T-19.1 [Front] Botón de auxilio de un solo toque (High) · T-19.2 [Back/Realtime] Notificación con alarma sonora prioritaria (High) · T-19.3 [Lógica] Registro de última sala del estudianteal disparar (Medium).
\item \textbf{HU-21 (Feedback caritas):} T-21.1 [Front] Modal de 3 caritas emocionales al terminar (Medium) · T-21.2 [Back] Registro analítico del estado de ánimo (Low).
\item \textbf{HU-22 (Plan B):} T-22.1 [BD] Configuración de tarea sustituta para imprevistos (High) · T-22.2 [Front] Botón de conmutación inmediata a Plan B (High) · T-22.3 [UX] Notificación tranquilizadora de cambio de plan (High).
\item \textbf{HU-23 (Bloqueo de tareas):} T-23.1 [Lógica] Jerarquía de prioridades que impide cancelar tareas críticas (Medium) · T-23.2 [Front] Función docente para suspender tareas con candado (Medium).
\item \textbf{HU-25 (Analítica semanal):} T-25.1 [Back] Consolidación de métricas de cumplimiento (Low) · T-25.2 [Admin] Informes gráficos de evolución de autonomía (Low).
\end{itemize}

\section{Actividad 7: Gestión del Backlog en JIRA Software Cloud}
\subsection{Estructura real configurada en JIRA}
El proyecto ha sido completamente inicializado y poblado en la plataforma oficial \hl{Atlassian JIRA Software Cloud}:

\begin{itemize}
\item \textbf{Nombre del proyecto:} \hl{GANTT-STERS}
\item \textbf{Clave del proyecto (Key):} \hl{SCRUM}
\item \textbf{Tipo de proyecto:} Scrum Board (Team-managed Software Development)
\item \textbf{Estado de carga en la herramienta:}
\begin{itemize}
\item \textbf{7 Épicas creadas:} \texttt{SCRUM-43} a \texttt{SCRUM-48} y \texttt{SCRUM-153} (Épica Técnica).
\item \textbf{26 Historias de Usuario cargadas}, con sus Story Points, prioridades, descripciones y criterios BDD.
\item \textbf{77 Subtareas técnicas de ingeniería} vinculadas dentro de cada historia.
\item \textbf{3 Sprints oficiales planificados} con sus fechas de inicio/fin y metas de entrega:
\begin{itemize}
\item \texttt{TUAG Sprint 1 (Iteración 1)}: 12 historias (9 funcionales + 3 HUT) · \textbf{34.0 Story Points}.
\item \texttt{TUAG Sprint 2 (Iteración 2)}: 7 historias · \textbf{29.0 Story Points}.
\item \texttt{TUAG Sprint 3 (Iteración 3)}: 7 historias · \textbf{23.0 Story Points}.
\end{itemize}
\end{itemize}
\end{itemize}

\subsection{Documento de acceso para el profesorado evaluador}
En cumplimiento estricto del apartado 5 del guion de la práctica (\textit{``Documento con información para acceder al proyecto definido en la herramienta JIRA''}), se aportan a continuación todos los cauces de acceso e inspección docente:

\begin{verde}{Datos de Acceso Oficial al Proyecto JIRA}
\textbf{Enlace directo al Tablero Scrum y Backlog:} \url{https://gantt-sters.atlassian.net/jira/software/projects/SCRUM/boards/1/backlog}\\
\textbf{Enlace al Tablero de Trabajo Activo:} \url{https://gantt-sters.atlassian.net/jira/software/projects/SCRUM/boards/1}\\
\textbf{Modalidad de acceso:}
\begin{enumerate}[leftmargin=1.6em,itemsep=2pt]
\item \textbf{Invitación nominal institucional:} Se ha remitido invitación formal con permisos de visualización y auditoría al correo institucional de los docentes de la asignatura (\texttt{@ugr.es}).
\end{enumerate}
\textbf{Permisos habilitados:} Inspección completa del Backlog, desglose de incidencias, criterios de aceptación, gráficos de evolución de sprint (\textit{Burndown Chart}) y panel de configuración de incidencias.
\end{verde}

\section{Conclusiones y estado de la entrega}
Con la redacción de este documento y la parametrización en tiempo real de la herramienta JIRA Cloud, el subgrupo \hl{Los Gantt-sters} concluye satisfactoriamente la totalidad de los requisitos de la \textbf{Práctica 2}:

\begin{enumerate}[leftmargin=1.6em,itemsep=2pt]
\item Se ha estructurado un Product Backlog integral de \textbf{26 Historias de Usuario} (23 funcionales y 3 técnicas) alineadas con las necesidades reales de los estudiantes y docentes del CEE Fundación Purísima Concepción.
\item El Product Owner (\textbf{Jaime Montávez Ruiz}) ha priorizado el catálogo optimizando el valor educativo y la resiliencia psicológica frente al cambio.
\item Se ha establecido una estimación matemática coherente mediante Planning Poker, vinculada a una velocidad realista de 31 Story Points por iteración (86 SP en total).
\item Se han definido exhaustivamente las 12 historias del primer ciclo con criterios BDD y 77 tareas técnicas para su ejecución inmediata.
\item El proyecto se encuentra 100\,\% operativo y auditable en JIRA Software Cloud.
\end{enumerate}

\end{document}
